\documentclass[12pt,letterpaper]{article}
\usepackage[utf8]{inputenc}
\usepackage[spanish]{babel}
\usepackage[margin=2.54cm]{geometry}
\usepackage{mathptmx} % Equivalente a Times New Roman
\usepackage{setspace}
\usepackage{enumitem}
\usepackage{hyperref}

\setlength{\parindent}{0pt}
\setlength{\parskip}{6pt}
\singlespacing

\begin{document}

\begin{center}
\textbf{\Large El Trío que Cuida los Datos}
\end{center}

\begin{center}
\textit{P.L. Alvarez} \\
\textit{7690 22 5839} \\
Universidad Mariano Gálvez \\
Seminario de Tecnologías de Información \\
\textit{palvarezg1@miumg.edu.gt}
\end{center}

\vspace{0.5cm}

\textbf{Resumen}

El objetivo de este trabajo es aclarar la diferencia entre tres términos que la gente suele usar como si fueran lo mismo: ciberseguridad, seguridad informática y seguridad de la información. Aunque los tres buscan proteger algo importante, cada uno tiene un enfoque distinto y confundirlos puede llevar a que una empresa invierta en la herramienta equivocada pensando que ya está protegida. Para lograr esta diferenciación se investigó y se tomó como base el modelo de los tres pilares de la seguridad de la información, que son confidencialidad, integridad y disponibilidad. Además del concepto del eslabón débil, que señala al ser humano como el punto más vulnerable de cualquier sistema de protección. Como resultado se logró explicar con ejemplos claros en qué se parecen y en qué se diferencian estos tres conceptos. Se concluye que la seguridad de la información es el concepto más amplio, que la seguridad informática se enfoca en proteger la infraestructura tecnológica, y que la ciberseguridad se centra específicamente en las amenazas que vienen del ciberespacio.

\textbf{Palabras clave:} ciberseguridad, seguridad informática, seguridad de la información, confidencialidad, integridad, disponibilidad.

\textbf{Desarrollo del tema}

Cuando alguien dice "necesito seguridad para mi empresa", en realidad puede estar hablando de tres cosas distintas y es importante saber cuál es cuál antes de invertir tiempo y dinero en protegerse. Los tres conceptos comparten un mismo objetivo general, que es evitar que algo importante se pierda, se dañe o caiga en manos equivocadas, pero cada uno ataca el problema desde un ángulo diferente.

La seguridad de la información es el concepto más amplio de los tres. No se enfoca únicamente en lo digital, sino en proteger la información sin importar el formato en el que esté, ya sea un archivo digital, un documento impreso o hasta una conversación. Este concepto se apoya en tres pilares fundamentales: la confidencialidad, que significa que solo las personas autorizadas pueden acceder a la información, la integridad, que significa que la información no se altera ni se modifica sin permiso y la disponibilidad, que significa que la información esté accesible cuando se necesite. Estos tres pilares forman lo que comúnmente se conoce como la triada de la información, y cualquier medida de protección que se implemente debería apuntar a cuidar estos tres elementos a la vez.

La seguridad informática, en cambio, es un concepto más específico dentro de la seguridad de la información. Se enfoca exclusivamente en proteger la infraestructura tecnológica, es decir, los equipos, los servidores, las redes y los sistemas donde vive esa información digital. Aquí es donde entran herramientas técnicas muy conocidas como los firewalls, los antivirus, las copias de respaldo en la nube y la seguridad en servidores. La seguridad informática se preocupa por amenazas como virus, troyanos, gusanos, spyware, backdoors, keyloggers y adwares, que son programas maliciosos diseñados para infiltrarse en los equipos y robar o dañar información. En pocas palabras, la seguridad informática protege el medio como son las computadoras y las redes, mientras que la seguridad de la información protege el contenidos es decir los datos en sí, sin importar dónde estén.

La ciberseguridad, por su parte, es un concepto que se enfoca específicamente en las amenazas que provienen del Internet y de los ataques cibernéticos en general. Mientras que la seguridad informática puede incluir amenazas físicas como que alguien robe una laptop o amenazas locales dentro de una red interna, la ciberseguridad se concentra en los ataques que llegan desde afuera a través de la red, como el phishing, el ransomware o los ataques dirigidos por hackers. Se podría decir que la ciberseguridad es como una rama especializada que nace de la seguridad informática, pero con el objetivo aen el mundo conectado a Internet.

Un concepto clave que atraviesa a los tres términos es el del eslabón débil de la seguridad. Por más firewalls, antivirus y protocolos que una empresa tenga instalados, el punto más vulnerable de cualquier sistema de seguridad suele ser el ser humano (capa 8). Un empleado que da clic en un enlace sospechoso, que usa una contraseña débil o que comparte información sin pensarlo puede tirar por la borda toda la inversión en herramientas técnicas. Por eso, cualquier estrategia de seguridad completa debe incluir no solo tecnología, sino también capacitación constante a las personas que usan esa tecnología, sirve mucho haciendo campañas de concientización.

Otra herramienta que conecta a los tres conceptos es la criptografía, que consiste en codificar la información para que solo las personas con el permiso y la capacidad de descifrarla puedan leerla. La criptografía se usa tanto para proteger la confidencialidad de la información (seguridad de la información) como para proteger las comunicaciones dentro de una red (seguridad informática) y para blindar las transacciones que viajan por Internet (ciberseguridad).

\textbf{Observaciones y comentarios}

Durante esta investigación se pudo notar que en la práctica muchas empresas y hasta algunos medios de comunicación usan estos tres términos como sinónimos, lo cual genera confusión a la hora de definir responsabilidades dentro de un equipo de trabajo. Por ejemplo, contratar a alguien que solo sabe de ciberseguridad no significa que la empresa esté protegida contra el robo físico de un servidor, porque eso corresponde más a la seguridad informática. También se observó que, sin importar cuál de los tres conceptos se aplique, el factor humano sigue siendo el eslabón más débil, lo cual demuestra que la tecnología por sí sola nunca es suficiente. Una limitación de este trabajo es que no se profundizó en casos reales de ataques donde se pueda ver claramente la diferencia entre los tres conceptos aplicados, lo cual sería un buen tema para ampliar en una investigación futura.

\textbf{Conclusiones}

\begin{enumerate}[label=\arabic*.]
\item La seguridad de la información es el concepto más amplio y se basa en proteger la confidencialidad, integridad y disponibilidad de los datos, sin importar el formato en el que estén.
\item La seguridad informática se enfoca específicamente en proteger la infraestructura tecnológica, como equipos, redes y servidores.
\item La ciberseguridad se concentra en las amenazas que vienen específicamente del ciberespacio e Internet.
\item El ser humano sigue siendo el eslabón más débil en cualquiera de los tres enfoques, por lo que la capacitación es tan importante como la tecnología.
\end{enumerate}

\textbf{Referencias}

Areitio, J. (2008). \textit{Seguridad de la información: Redes, informática y sistemas de información}. Cengage Learning.

International Organization for Standardization. (2022). \textit{ISO/IEC 27001:2022 Information security management systems — Requirements}.

National Institute of Standards and Technology. (2018). \textit{Framework for improving critical infrastructure cybersecurity} (versión 1.1). NIST.

\vspace{0.5cm}
\noindent \small Repositorio: https://github.com/PLAlvarez11/seminario-encuesta.git

\end{document}